\chapter*{The eight-part series}
\addcontentsline{toc}{chapter}{The eight-part series}
The series now has eight parts. Its first two form a continuous course:
understand the physical account, derive it carefully, inspect how it can be
implemented, and learn how evidence limits a claim. Later volumes deepen
specific research questions; their position in the list is not execution
permission or a guarantee of a positive result.

\begin{longtable}{P{.10\linewidth}P{.66\linewidth}P{.12\linewidth}}
\toprule New & Responsibility & Former\\\midrule
I & Physical intuition, motivation and the EBU idea & I\\
II & Mathematics, implementation and evidence & II+III\\
III & Measurement, evidence and falsification & IV\\
IV & Homeostasis, recovery and stochastic dynamics & V\\
V & Multiple actions, interaction and local structure & VI\\
VI & Across distance & VII\\
VII & Coordination, memory and system dynamics & VIII\\
VIII & Institutions and an action-accounted economy & IX\\\bottomrule
\end{longtable}

Part III's methods and Part V's simultaneous foundations support the
stochastic questions in Part IV. Part V supports the route and boundary
work in Part VI. Parts IV--VI inform Part VII's explicit coordination
policies. Part VIII depends on surviving scientific and social evidence,
including adverse results, and on normative arguments that no accounting
identity can supply.

The introductory group treatment in this Part II is sufficient to explain
exact finite accounting and its boundaries. The later Part V has a different
job: deeper interaction diagnostics, admissible subset structures, topology,
canonical motifs and certified compression. Likewise, this volume's evidence
chapters teach how to read the existing record; Part III develops measurement
and falsification methods in greater depth. These are deliberate levels of
treatment, not reasons to repeat the same chapter in two places.

Historical references keep the numbering printed in their original edition.
The renumbering changes the forward editorial map only. It neither rewrites
a theorem citation nor changes an experimental programme's stage identifiers.
